% Throwaway smoke test for assets/preamble.srm.tex — run BEFORE spending anything on a
% conversion.  It exercises every decision the patch makes and then the .aux is read back
% to see what each \label actually resolved to.
%
%   cd .texify_work && pdflatex -interaction=nonstopmode smoke.tex && grep newlabel smoke.aux
%
% What it must show:
%   q:1 -> 1 ... q:75 -> 75          the argument, not an auto-increment
%   s:7 with NO "Key:" printed        the one solution that has none
%   choices lettered (A)..(E)         and a four-item statements list reaching IV
%   the 114pt figure NOT blown up, the 495pt figure shrunk to the measure
% texify — shared master preamble for JOURNAL ARTICLES.
%
% This is the committed base. `--master` inlines it verbatim, then \input's the
% per-paper patch assets/preamble.<paper>.tex, then opens the document. Nothing here
% is scraped from the source; edit it by hand.
%
% CLASS IS `article`, ON PURPOSE. The textbook pipeline this was forked from used
% `report`, where \thesection expands to \thechapter.\arabic{section}; with no
% \chapter ever issued the chapter counter stays 0 and a paper's "1. Introduction"
% prints as "0.1". That repo worked around it with three \counterwithout calls.
% `article` has no chapter at all, so the trap cannot occur -- and \chapter is
% undefined, which turns a converter that emits one into a hard compile error
% instead of a silent mis-numbering.
%
% NUMBERING IS FLAT HERE (equations 1,2,3...; theorems 1,2,3...). Papers differ:
% many number within sections ((2.1), Theorem 3.4). That is a per-paper decision, so
% the \counterwithin calls live in assets/preamble.<paper>.tex, not here.

\documentclass[11pt]{article}
\usepackage[margin=25mm]{geometry}

\usepackage{amsmath,amsthm,amssymb,amsfonts}   % amssymb: \leqslant \geqslant \rightleftarrows
\usepackage{mathtools}
\usepackage{booktabs,array}
\usepackage[shortlabels]{enumitem}
\usepackage{graphicx}
\usepackage{xcolor}
\graphicspath{{src/}{./}}

% Numeric [n] citations -- the dominant style in mathematics journals (prose names the
% author, then a bracketed number: "Ross [2, p. 67]"). A paper that cites author-year
% instead flips this at run time from its own preamble patch with
%     \setcitestyle{authoryear,round,aysep={},yysep={;},notesep={, }}
% -- package options cannot be changed after loading, so \setcitestyle is the only way
% to do it without editing this file. natbib loads BEFORE hyperref so cites hyperlink.
\usepackage[numbers,sort&compress]{natbib}

\usepackage{hyperref}
\hypersetup{colorlinks=true,linkcolor=blue!60!black,urlcolor=blue!70!black,
            citecolor=blue!60!black}

% -------------------------------------------------------------- theorem-likes --
% A generic paper set. `theorem` owns the counter and lemma/proposition/corollary/
% definition share it -- the usual mathematics-journal convention (Theorem 1, Lemma 2,
% Theorem 3, ...). A paper that numbers them independently, or that needs an
% environment not listed here, overrides in its own assets/preamble.<paper>.tex --
% never by editing this file. (Splitting an environment back off a shared amsthm
% counter is fiddly: see the \let\the<env>\relax note in assets/preamble.white.tex.)
\newtheorem{theorem}{Theorem}
\newtheorem{lemma}[theorem]{Lemma}
\newtheorem{proposition}[theorem]{Proposition}
\newtheorem{corollary}[theorem]{Corollary}

\theoremstyle{definition}
\newtheorem{definition}[theorem]{Definition}
\newtheorem{example}{Example}

\theoremstyle{remark}
\newtheorem*{remark}{Remark}

\theoremstyle{plain}

% ----------------------------------------------------------------- operators --
% House-fixed glyphs for probability/statistics. The conversion contract mandates these
% over \mathbb{P} / \mathbb{E} / \Pr / \mathbb{1}, so notation stays uniform across
% papers even when the sources disagree. Unused by a paper that has no such operators.
\usepackage{dsfont}
\DeclareMathOperator\prb{\mathsf{P}}
\DeclareMathOperator\expc{\mathsf{E}}
\DeclareMathOperator\var{var}
\DeclareMathOperator\cov{cov}
\DeclareMathOperator\cor{cor}
\DeclareMathOperator\indc{\mathds{1}}
\DeclareMathOperator\tr{tr}
\DeclareMathOperator\diag{diag}
\DeclareMathOperator\rk{rank}

% Differentials and Euler's number roman: \dd x, \e^{x}.
\newcommand{\dd}{\mathrm{d}}
\newcommand{\e}{\mathrm{e}}

\setlength{\jot}{4pt}

% srm — Society of Actuaries, "Exam SRM — Statistics for Risk Modeling:
% Sample Questions and Solutions" (February 2026 revision, 75 pp.)
%
% Per-document overrides. The master \input's this AFTER the shared assets/preamble.tex,
% so everything here layers on top of the committed base without touching it.
%
% HAND-WRITTEN from the page images (all 75 rendered at 150 dpi and read; the phi on p56
% and the MT-Extra ellipses re-cropped at 600 dpi).
%
% ⚠ THIS SOURCE IS NOT A PAPER AND NOT A BOOK. It is a flat list of 75 multiple-choice
% questions and their 75 solutions, and it has NO sections, NO theorems, NO numbered
% equations and NO reference list. Almost nothing in the base preamble's apparatus applies;
% what it needs instead is the four environments below. The base is still the right base —
% `article`, amsmath, graphicx, enumitem, booktabs — so this patch only ADDS.

% =============================================================== the four environments ==
%
% The deliverable interleaves the source's two halves: each question is immediately
% followed by its own solution. `texify_qa_merge.py` does the pairing, and it can only do
% that if each item is delimited and carries its printed number, hence the mandatory
% number argument on both `question` and `solution`.
%
% ⚠ THE NUMBER IS AN ARGUMENT, NOT AN AUTO-INCREMENT. 1-75 looks contiguous but is not
% self-checking: four numbers (17, 28, 47, 65) were withdrawn in August 2023 and the source
% still prints them, as the bare word DELETED. An auto-incrementing counter would renumber
% the 71 live questions the moment the merge dropped a deleted one, and nothing downstream
% would notice — the numbers would still run 1..71 without a gap. Taking the number from the
% page makes a mis-paired item a compile-visible wrong number instead.
%
% ⚠ AND THE COUNTER IS STILL STEPPED, via \refstepcounter, so that \label records what
% LaTeX actually set. That is what lets verify.py read `\newlabel{q:33}{{33}...}` back out
% of the .aux and assert the printed number equals the number the page shows — the one
% check that catches a numbering shift from a cause nobody thought to test for.

% ⚠ \makeatletter IS LOAD-BEARING. This file is \input from the master's preamble, and an
% \input'ed .tex is NOT a .sty: `@` keeps catcode 12 there, so the private \srm@... names
% below would parse as \srm followed by the letters "@itemlist" and the first \newcommand
% would die with "Missing \begin{document}" — an error that names neither the macro nor the
% file. (It did, on the first smoke run.)
\makeatletter

\newcounter{srmq}
\newcounter{srms}

% The question's geometry: the number hangs at the left margin and the body is indented
% past it, as the source sets it. 2.2em holds "75." with room to spare.
\newcommand{\srm@itemlist}{%
  \setlength{\labelwidth}{2.2em}%
  \setlength{\labelsep}{0.6em}%
  \setlength{\leftmargin}{2.8em}%
  \setlength{\itemindent}{0pt}%
  \setlength{\listparindent}{0pt}%
  \setlength{\parsep}{0.4\baselineskip}%
  \setlength{\itemsep}{0pt}%
  \setlength{\topsep}{0pt}%
  \setlength{\partopsep}{0pt}%
}

% The solution's geometry: NO label box, and the same 2.8em left edge as the question's
% BODY, so a solution reads as subordinate to the question it now sits under.
%
% ⚠ THE HEAD MUST NOT BE A LIST LABEL. The obvious \item[\textbf{Solution.}] puts the word
% in a \labelwidth-wide box, and LaTeX right-aligns a label in that box and lets an oversize
% one hang out to the LEFT. "Solution." is about twice 2.2em, so it hung 2em past the list's
% own margin and every solution head printed further left than the question number above it
% — the visual hierarchy exactly inverted. Widening \labelwidth instead would have indented
% every solution body 5em. So: an empty label, and the head run in as ordinary bold text.
\newcommand{\srm@sollist}{%
  \setlength{\labelwidth}{0pt}%
  \setlength{\labelsep}{0pt}%
  \setlength{\leftmargin}{2.8em}%
  \setlength{\itemindent}{0pt}%
  \setlength{\listparindent}{0pt}%
  \setlength{\parsep}{0.4\baselineskip}%
  \setlength{\itemsep}{0pt}%
  \setlength{\topsep}{0pt}%
  \setlength{\partopsep}{0pt}%
}

% \begin{question}{33} ... \end{question}
\newenvironment{question}[1]{%
  \par\addvspace{1.4\baselineskip}%
  \setcounter{srmq}{\numexpr#1-1\relax}\refstepcounter{srmq}%
  \begin{list}{}{\srm@itemlist}%
  \item[\textbf{\thesrmq.}]\label{q:#1}\ignorespaces
}{%
  \end{list}%
}

% \begin{solution}{33}{E} ... \end{solution}
%
% ⚠ THE KEY ARGUMENT MAY BE EMPTY, and for exactly one solution it must be: solution 7
% (p58) carries no "Key:" line at all — it is four lines of prose and nothing else. A
% one-argument environment would have forced a key to be invented for it.
%
% The three printed spellings "Key: E", "Key E" (solution 13) and "Key:D" (solutions 66 and
% 75) are NORMALISED to one here, because the head is re-typeset from the letter rather than
% transcribed. That is the one deliberate departure from reproduce-what-is-printed in this
% conversion; it is recorded in README.md and asserted by verify.py, which checks all 71
% letters against the page.
\newcommand{\srm@key}[1]{\if\relax\detokenize{#1}\relax\else\ \textbf{Key:~#1}\fi}
\newenvironment{solution}[2]{%
  \par\addvspace{0.7\baselineskip}%
  \setcounter{srms}{\numexpr#1-1\relax}\refstepcounter{srms}%
  \begin{list}{}{\srm@sollist}%
  \item\label{s:#1}\textbf{Solution.}\srm@key{#2}\quad\ignorespaces
}{%
  \end{list}\addvspace{0.3\baselineskip}%
}

% The five answer choices, printed "(A)" ... "(E)". \Alph* is the letter, so a dropped
% choice re-letters the ones after it — which verify.py catches by counting five per
% question and by checking the rendered letters against the .aux.
\newlist{choices}{enumerate}{1}
\setlist[choices]{label=(\Alph*), ref=(\Alph*), align=left,
                  leftmargin=2.8em, labelwidth=2.2em, labelsep=0.6em,
                  topsep=0.6\baselineskip, itemsep=0.15\baselineskip, parsep=0pt}

% The roman-numbered statement list a question offers for "Determine which of the following
% statements is/are true": I., II., III., and on questions 5, 6 and 74 also IV.
\newlist{statements}{enumerate}{1}
\setlist[statements]{label=\Roman*., ref=\Roman*, align=left,
                     leftmargin=3.2em, labelwidth=2.6em, labelsep=0.6em,
                     topsep=0.6\baselineskip, itemsep=0.15\baselineskip, parsep=0pt}

% The lower-case roman "given" list: i), ii), iii), iv). Note the CLOSING PAREN ONLY — the
% source prints "i)", never "(i)", and never a full stop.
\newlist{given}{enumerate}{1}
\setlist[given]{label=\roman*), ref=\roman*, align=left,
                leftmargin=2.8em, labelwidth=2.2em, labelsep=0.6em,
                topsep=0.4\baselineskip, itemsep=0.15\baselineskip, parsep=0pt}

% ========================================================================= paragraphs ==
%
% The source is a Word document set BLOCK STYLE throughout: not one paragraph on any of its
% 75 pages is first-line indented, and paragraphs are separated by vertical space instead.
% `article`'s default is the opposite, so without this the front matter's 18 update lines —
% each its own paragraph — would print as an indented staircase, and every solution's second
% paragraph would step in under the first.
%
% A modest fixed \parskip with a little stretch: enough to separate, not so springy that the
% lists (whose \parsep and \topsep are set explicitly above) start to wobble.
\setlength{\parindent}{0pt}
\setlength{\parskip}{0.5\baselineskip plus 0.15\baselineskip minus 0.1\baselineskip}

% ========================================================================== furniture ==

% The word the source prints, alone, for each of the four withdrawn numbers. Upright roman
% capitals, exactly as on the page — NOT \textsc, which would set small capitals.
\newcommand{\deleted}{DELETED}

% The two half-titles, "QUESTIONS" (p2) and "SOLUTIONS" (p57). The converter emits one at
% the head of its division; `texify_qa_merge.py` KEEPS them in `--order original` and DROPS
% them in the default interleaved order, where a heading announcing a half that no longer
% exists would be a lie.
\newcommand{\srmpart}[1]{%
  \par\addvspace{1.2\baselineskip}%
  {\centering\bfseries #1\par}%
  \addvspace{0.4\baselineskip}%
}

% =========================================================================== figures ==
%
% 17 figures, cropped out of the source by texify_figs_srm.py into assets/figs/srm_<key>.pdf
% and placed with \srmfig{<key>}. Keys are per QUESTION (q35, q48c), never per division —
% see SKILL.srm.md for the full table.
%
% ⚠ THIS IS NOT THE PIPELINE'S USUAL `<book>_<NN>_figs.pdf` + `page=<k>` CONVENTION, and
% deliberately so. That convention indexes a figure by its position within a division's
% figure file, which here would mean the answer-choice picture of question 48 is "page 8 of
% srm_02_figs.pdf" — a number that silently changes meaning the moment a division boundary
% moves, with no way for a reader or a check to notice. A key naming the question cannot
% drift.
%
% Natural size when it fits the measure, shrunk to the measure when it does not: two of the
% seventeen (q39 at 474pt, q57 at 495pt) are wider than the 470pt text block, and the three
% panels of question 26 are 114pt and must NOT be blown up to fill it.
%
% ⚠ AND A FIGURE IS OFTEN THE WHOLE OF A LIST ITEM — question 48's five answer choices are
% each nothing but a picture, and question 26's three roman items likewise. A list label is
% not typeset by \item; it is parked and inserted by \everypar when the item's first
% PARAGRAPH starts. If that first paragraph is the centred graphic, the label is prepended to
% it and lands on the picture's BASELINE — i.e. at the bottom-left of a 140pt-tall panel,
% where "(A)" reads as if it belonged to the choice below. (Observed: all five of question
% 48's letters.)
%
% \if@inlabel is exactly "a label is parked and waiting", so open and close an empty
% paragraph to let it out FIRST, and only there. Testing \ifhmode instead cannot work: an
% item's start and an ordinary between-paragraphs position are both vertical mode, and
% unconditionally emitting the empty paragraph would pad every one of the other nine figures
% with a blank line.
\newsavebox{\srm@figbox}
\newcommand{\srmfig}[1]{%
  \if@inlabel\leavevmode\par\fi
  \par\addvspace{0.5\baselineskip}%
  \begin{center}%
    \sbox{\srm@figbox}{\includegraphics{assets/figs/srm_#1.pdf}}%
    \ifdim\wd\srm@figbox>\linewidth
      \resizebox{\linewidth}{!}{\usebox{\srm@figbox}}%
    \else
      \usebox{\srm@figbox}%
    \fi
  \end{center}%
  \addvspace{0.5\baselineskip}%
}

% ========================================================================= operators ==
%
% The base declares \expc as \mathsf{E} and \var as an upright "var", the fleet's house
% style. This source prints ITALIC "E(...)" and italic "Var(...)", and it is a study aid for
% candidates who will meet the SOA's notation in the exam room, so the macros are re-pointed
% at the source's own glyphs. The macro NAMES stay the base's, so the house rule "never
% \mathbb{E}, never a bare E" still binds and the notation is uniform across all 75 items.
%
% \DeclareMathOperator is \newcommand underneath and refuses to redefine; \let ...\relax
% first is the documented way round, and is the same trick assets/preamble.tex's siblings
% use to free an amsthm counter.
\let\expc\relax  \DeclareMathOperator{\expc}{\mathit{E}}
\let\var\relax   \DeclareMathOperator{\var}{\mathit{Var}}
\let\cov\relax   \DeclareMathOperator{\cov}{\mathit{Cov}}
\let\cor\relax   \DeclareMathOperator{\cor}{\mathit{Corr}}

% ====================================================================== bibliography ==
%
% NOTHING TO DO, AND NOTHING TO ADD. This source has no reference list. It names its
% textbooks in running prose — "See Page 242 of Regression Modeling with Actuarial and
% Financial Applications", "see page 307 of Frees", "Section 8.1 of An Introduction to
% Statistical Learning" — with the title in italics and the page as ordinary words. Those
% are TEXT, not citations: there is no \cite anywhere in the body, natbib is loaded by the
% base but never used, and the master emits no \bibliography. A .bib invented to hold them
% would print a reference list the source does not have.

\makeatother

\begin{document}

\srmpart{QUESTIONS}

\begin{question}{1}
You are given the following four pairs of observations:
\[ x_1 = (-1,0), \quad x_2 = (1,1), \quad x_3 = (2,-1), \text{ and } x_4 = (5,10). \]
Calculate the intercluster dissimilarity between $\{x_1, x_2\}$ and $\{x_4\}$.
\begin{choices}
\item 2.2
\item 3.2
\item 9.9
\item 10.8
\item 11.7
\end{choices}
\end{question}

\begin{question}{5}
Consider the following statements:
\begin{statements}
\item Principal Component Analysis (PCA) provide low-dimensional linear surfaces.
\item The first principal component is the line in p-dimensional space.
\item PCA finds a low dimension representation of a dataset.
\item PCA serves as a tool for data visualization.
\end{statements}
Determine which of the statements are correct.
\begin{choices}
\item Statements I, II, and III only
\item Statements I, II, and IV only
\item Statements I, III, and IV only
\item Statements II, III, and IV only
\item Statements I, II, III, and IV are all correct
\end{choices}
\end{question}

\begin{question}{17}
\deleted
\end{question}

\begin{question}{26}
Each picture below represents a two-dimensional space.
\srmfig{q26a}
\srmfig{q26b}
\end{question}

\begin{question}{21}
A random walk is expressed as
\[ y_t = y_{t-1} + c_t \quad \text{for } t = 1,2,\ldots \]
where $\expc(c_t) = \mu_c$ and $\var(c_t) = \sigma_c^2$, $t = 1,2,\ldots$, and
$\phi\mu^p \geqslant 0$, $\lambda \leqslant 1$.
\begin{given}
\item The random walk model $y_t = y_0 + c_1 + c_2 + \cdots + c_t$.
\item The following nine observed values of $c_t$:
\begin{center}\begin{tabular}{|c|c|c|c|}\hline
$t$ & 11 & 12 & 13 \\ \hline
$c_t$ & 2 & 3 & 5 \\ \hline
\end{tabular}\end{center}
\end{given}
\begin{choices}
\item 1
\item 2
\item 3
\item 8
\item 18
\end{choices}
\end{question}

\begin{question}{57}
The following plot of the corresponding regression tree:
\srmfig{q57}
\end{question}

\begin{question}{75}
Determine which of the following statements about the Tweedie distribution is NOT true.
\begin{choices}
\item Its variance is $\phi\mu^{p}$.
\item The probability mass at zero is $e^{-\lambda}$.
\item It is a member of the linear exponential family.
\item The Poisson distribution is a special case.
\item This Tweedie distribution is a Poisson sum of gamma random variables.
\end{choices}
\end{question}

\begin{question}{48}
The following tree was constructed using recursive binary splitting.
\srmfig{q48}
Determine which of the following plots represents this tree.
\begin{choices}
\item \srmfig{q48a}
\item \srmfig{q48b}
\item \srmfig{q48c}
\item \srmfig{q48d}
\item \srmfig{q48e}
\end{choices}
\end{question}

\begin{question}{26}
Determine which space can be modeled with no error by a classification tree.
\begin{statements}
\item \srmfig{q26a}
\item \srmfig{q26b}
\item \srmfig{q26c}
\end{statements}
\begin{choices}
\item I only
\item II only
\item III only
\item I, II, and III
\item The correct answer is not given by (A), (B), (C), or (D).
\end{choices}
\end{question}

\srmpart{SOLUTIONS}

\begin{solution}{1}{E}
For complete linkage, the dissimilarity measure used is the maximum, which is 11.66.
\end{solution}

\begin{solution}{7}{}
The intent is to model a binary outcome, thus a classification model is desired.
This solution carries NO key line in the source, and none must be printed here.
\end{solution}

\begin{solution}{17}{}
\deleted
\end{solution}

\begin{solution}{75}{D}
D is not correct. As $p$ approaches 1, this distribution approaches the Poisson
distribution.
\end{solution}

\end{document}
